\chapter{Literature Review}
\label{ch:literature}

\section{Introduction}
This chapter reviews the ideas behind the project's retrieval pipeline and
explains the choices made in applying them. The literature supplies model
architectures and evaluation concepts; the project's own experiments, reported
in Chapter~\ref{ch:evaluation}, supply evidence about their use on the available
clothing gallery. Results reported by a model's authors are not substituted for
measurements of this application.

\section{Fashion Retrieval and the Meaning of Relevance}
\subsection{Category, Resemblance, and Identity}
Fashion search can answer several different questions. Category retrieval asks
whether a result is, for example, a dress rather than a shirt. Resemblance asks
whether its shape, colour, material appearance, or pattern is useful to the
searcher. Identity retrieval asks whether two photographs show the same product.
These definitions can disagree: a plain grey dress and a patterned dress may
share a category while differing visibly, and two alternate views may show the
same product despite changes in pose and lighting.

DeepFashion made clothing recognition and retrieval more explicit through rich
annotations and image correspondences in different capture settings
\cite{Liu2016DeepFashion}. DeepFashion2 further addresses detection, landmarks,
segmentation, and clothing retrieval in a common benchmark
\cite{Ge2019DeepFashion2}. These works demonstrate why a retrieval study needs
labels appropriate to its intended claim. An image category alone cannot provide
the ground truth for exact product matching.

\subsection{Implications for the Present Dataset}
The project uses the ten-class clothing subset published by Grigorev
\cite{GrigorevClothingDataset}. It is practical for local experimentation and
category-level comparisons. However, its supplied image categories do not
establish merchant SKU identities. The application's file-derived identifiers
serve indexing and display; they do not turn category labels into identity
annotations. The study therefore reports category retrieval as its principal
quantitative evidence and treats visually reviewed alternate views as a separate,
limited identity pilot.

This distinction also explains a common user observation. A search may return
an appropriate garment type while ranking a less convincing colour above a
closer-looking alternative. That is a potential fine-grained ranking problem,
but a category-only metric may count both results as equally relevant. A
future resemblance study requires independently judged candidate relevance,
not an assumption that every same-category candidate is equally useful.

\section{Pretrained Visual Representations}
\subsection{Convolutional Baselines}
Residual networks introduce shortcut connections that support the optimisation
of deep visual models~\cite{He2015ResNet}. ResNet50 provides a familiar
convolutional baseline from which features can be extracted without training
a clothing-specific network. EfficientNet studies the coordinated scaling of
network depth, width, and resolution~\cite{Tan2019EfficientNet}; the B3 variant
is another baseline in the project's saved comparison. Their inclusion tests
whether a general convolutional representation is sufficient for this gallery.

Features useful for classification need not order visually similar catalogue
items in the way a shopper expects. In this project, the baselines are compared
as fixed feature extractors under a common retrieval task. Their local results
are evidence about those configurations and that dataset, rather than a general
ranking of all convolutional and transformer architectures.

\subsection{Vision-Language and Self-Supervised Encoders}
CLIP learns transferable visual representations through paired image and text
supervision~\cite{Radford2021CLIP}. Its image encoder can be used for image-to-image
retrieval even when no text query is entered. DINOv2 instead develops general
visual features through self-supervised learning and curated visual data
\cite{Oquab2023DINOv2}. These two families offer useful contrasting baselines:
semantic image-text training and visual self-supervision.

SigLIP2 extends a vision-language representation-learning recipe with additional
objectives and data curation~\cite{Tschannen2025SigLIP2}. The project adopts
the published SigLIP2 Base patch16-224 checkpoint with its documented
image processor~\cite{GoogleSigLIP2Card}. Its image branch produces the
768-dimensional vectors used by the application. The original checkpoint is
loaded locally and kept fixed; the project does not train a new vision-language
model or use text matching in the implemented user search.

The saved local validation comparison favours SigLIP2 among the five tested
configurations. This selection is supported by the project's retained results,
not solely by the model paper. It also keeps the deployed pipeline focused on
one retrieval encoder. The four other encoders remain experimental baselines
and are not required to run the presentation application.

\section{Interactive Segmentation}
\subsection{From a Full Photograph to a Target Garment}
A global representation can mix information from several garments and the
background. User-guided segmentation addresses the target-selection ambiguity
before retrieval. SAM2 is a promptable segmentation model for images and video
\cite{Ravi2024SAM2}. The official release also provides the SAM2.1 checkpoints
\cite{MetaSAM2Repository}; this project uses the Hiera Base Plus checkpoint
distributed through the model repository~\cite{MetaSAM21Card}.

Only still-image point prompting is exposed in the application. Positive points
identify garment regions and negative points exclude unwanted regions. Additional
points can improve an incomplete selection, such as a mask that initially
captures a decoration rather than the whole garment. The user inspects the
overlay and accepts the cutout before searching. This makes the user part of
the selection process and gives a practical recovery path when one click is
insufficient.

\subsection{Selection Quality and Retrieval Quality}
The model returns a candidate-mask quality score, while retrieval returns a
vector similarity score. Neither is an observed measurement of pixel-level
segmentation accuracy or exact-product identity. Measuring mask accuracy would
require independent reference masks and an overlap measure such as IoU. The
project does not possess that annotated evaluation set and therefore does not
publish a mask-IoU accuracy result.

Manual rectangle, polygon, and freehand selection provide alternatives when the
model's mask is unsuitable. The manual region and the accepted SAM2 cutout both
enter the same retrieval path. This design separates the question of which pixels
the user selects from the question of how the selected image is ranked.

\section{Vector Similarity Search}
\subsection{Normalised Embeddings and Cosine Similarity}
Let $f(I)\in\mathbb{R}^{d}$ be an image representation and let
$\mathbf{z}=f(I)/\lVert f(I)\rVert_2$ be its normalised vector. For a query
$\mathbf{z}_q$ and a catalogue vector $\mathbf{z}_i$, the scoring rule is
\begin{equation}
 s(q,i)=\mathbf{z}_q^{\mathsf T}\mathbf{z}_i
 =\frac{f(I_q)^{\mathsf T} f(I_i)}{\lVert f(I_q)\rVert_2\lVert f(I_i)\rVert_2}.
 \label{eq:cosine}
\end{equation}
Normalisation makes inner-product ranking equivalent to cosine ranking. FAISS
provides both exact and approximate vector-search structures
\cite{Johnson2017FAISS,FaissIndexes}. With 3,068 vectors, the project uses exact
\texttt{IndexFlatIP}; it does not need an approximate-neighbour index to demonstrate
the present workflow. Neural inference may use CUDA, but the installed FAISS
index runs on the CPU.

\subsection{Category Filtering and Grouped Results}
Ranking can be constrained by a predicted or user-selected clothing category.
The original gallery uses normalised category centroids derived from its own
embeddings. Filtering can remove a strongly ranked item from the wrong category,
but an incorrect category prediction can also remove a useful answer. The
interface therefore provides automatic, explicit-category, and unrestricted
search modes.

The managed catalogue has a different data structure: several images can
belong to one product. Its ranking retains the best matching photograph and
returns each product once. A separate, explicitly enabled reviewed-group option
can merge previously reviewed identities in the original gallery. Neither
mechanism infers that arbitrary similar images show the same commercial product.

\section{Evaluation Concepts}
Precision at a rank cutoff, average precision, and reciprocal rank describe
different properties of a ranked list~\cite{Manning2008IR}. Precision asks how
many early results are relevant. Reciprocal rank rewards an early first relevant
result. Average precision combines precision at relevant ranks. Their meaning
depends on the relevance labels used; the same formula can measure category
relevance or a verified identity relationship. Chapter~\ref{ch:methodology}
defines the exact truncated formulas used by the project.

A valid comparison also requires control over the gallery, query images,
preprocessing, and selection prompts. Reusing a cached embedding file can be
useful for repeatability, but cache existence alone does not establish that
the file matches the current model and code. For this reason, the report
distinguishes retained historical scores from application-aligned comparisons
with recorded artifact hashes.

\section{Application and Data-Management Foundations}
FastAPI provides a Python API framework with typed request handling and generated
API documentation~\cite{FastAPIDocumentation}. It is used to connect uploaded
images and point prompts to local inference services. A browser frontend keeps
the workflow accessible without packaging the neural models into client-side
JavaScript. SQLite provides transactional local persistence
\cite{SQLiteTransactions}, fitting the current single-computer product manager.

Local administration uses password hashing through \texttt{argon2-cffi}
\cite{Argon2Parameters}, revocable sessions, and server-side access checks.
These components support the implemented operating scope. They are not evidence
that the application has undergone a production security audit or can already
support independent seller organisations on a public service.

\section{Comparison and Project Rationale}
\begin{table}[H]
\centering\caption{Research foundations and their role in this project}
\label{tab:related-work}
\begin{ReportTable}
\begin{tabularx}{\linewidth}{@{}L{0.22\linewidth}Y Y@{}}
\toprule\textbf{Work or component}&\textbf{Relevant contribution}&\textbf{Project use or boundary}\\\midrule
DeepFashion and DeepFashion2 & Fashion annotations and cross-view retrieval tasks & Inform relevance definitions; not the deployed dataset \\ \addlinespace
ResNet50 and EfficientNet-B3 & General convolutional representations & Historical retrieval baselines \\ \addlinespace
CLIP and DINOv2 & Transferable semantic or self-supervised features & Alternative pretrained encoder baselines \\ \addlinespace
SigLIP2 & Pretrained image representation & Adopted fixed image encoder \\ \addlinespace
SAM2 / SAM2.1 & Promptable visual segmentation & Interactive still-image garment cutout \\ \addlinespace
FAISS & Exact and approximate vector indexing & Exact CPU inner-product search \\
\bottomrule
\end{tabularx}
\end{ReportTable}
\end{table}

The project addresses an integration problem: communicate the desired garment,
retain an inspectable query, rank a bounded gallery, and maintain a searchable
collection through a local web interface. The contribution is the working
combination, its safeguards and evaluation, rather than a claim of a new
foundation-model architecture. Table~\ref{tab:related-work} relates each research
foundation to that role.

\section{Summary}
The literature supports separating visual representation, target selection,
ranking, and relevance evaluation. SigLIP2, SAM2.1, and FAISS are selected for
these complementary responsibilities. The next chapter specifies the data,
protocols, and metric conventions used to test their combination.
